\section{Introducción}
En el desarrollo de software, garantizar la calidad del producto final es un aspecto fundamental para satisfacer las necesidades de los usuarios y asegurar su correcto funcionamiento en diferentes entornos. Se han propuesto diversos modelos de evaluación, entre los cuales, el modelo de calidad de Boehm, organiza características del software en diferentes niveles o capas, permitiendo analizar tanto aspectos generales como detalles técnicos y específicos.\\ \\
Este trabajo presenta una descripción estructurada de dichas características, clasificándolas en niveles de alto nivel, intermedio y primitivo, con el objetivo de comprender cómo cada una contribuye a la calidad global del software. Además, se analiza la relación entre estas capas y su importancia en la evaluación del producto.
